% Consolidated 35-wave-properties from 620defd2149c8dcfd52d29c147aa4708c27ab8f3; original TeX fragments preserved.

% === BEGIN 35-wave-properties-049-065.tex blob ab3f96ba9c795b8c98b5a854d1f29a9c77d79557 ===
% Irodov Part VI §6.2 original problems 49--65.
% Source: supplied Arihant PDF pp.247--248, printed pp.238--239.
% Answers PDF pp.349--350. Every TeX derivation has an excheck clause.
\wavequantproblem{49}{238}
动能均为 $T=100\,\mathrm{eV}$ 的电子、
质子和铀原子，分别求其德布罗意波长。
\exsolution
三个粒子均在非相对论速度区间，
$p=\sqrt{2mT}$，所以
\[
 \boxed{\lambda_{\rm dB}=\frac h{\sqrt{2mT}}}.
\]
代
$m_e=9.109\times10^{-31}\,\mathrm{kg}$、
$m_p=1.673\times10^{-27}\,\mathrm{kg}$、
$m_{\rm U}\simeq238u$，得到
\[
 \boxed{\lambda_e\simeq123\,\mathrm{pm}},\qquad
 \boxed{\lambda_p\simeq2.86\,\mathrm{pm}},\qquad
 \boxed{\lambda_{\rm U}\simeq0.186\,\mathrm{pm}}.
\]
\excheck 相同动能下 $\lambda\propto m^{-1/2}$；
最重的铀原子波长最短。

\wavequantproblem{50}{238}
一个电子的德布罗意波长从
$100\,\mathrm{pm}$ 降到 $50\,\mathrm{pm}$，
问需要增加多少动能？
\exsolution
在这一能区电子可采用非相对论关系
\[
 T(\lambda)=\frac{h^2}{2m_e\lambda^2}.
\]
$\lambda_2=\lambda_1/2$，所以
$T_2=4T_1$，新增动能为
\[
 \boxed{\Delta T=T_2-T_1=
 \frac{3h^2}{2m_e(100\,\mathrm{pm})^2}
 \simeq0.45\,\mathrm{keV}}.
\]
\excheck 波长减半使动量加倍、动能增至4倍，
但\emph{增加}的是初动能的3倍。

\wavequantproblem{51}{238}
动能为 $T=25\,\mathrm{eV}$ 的中子撞击
起初静止的氘核。
求两粒子在质心参考系内各自的德布罗意波长。
\exsolution
实验室系中中子速度
$v_n=\sqrt{2T/m_n}$。
质心速度
$V_{\rm CM}=m_nv_n/(m_n+m_d)$。
中子质心动量模
\[
 p_n^*=m_n(v_n-V_{\rm CM})
 =\frac{m_nm_d}{m_n+m_d}v_n=\mu v_n.
\]
氘核的质心动量与它等大反向，
所以两者波长相同：
\[
 \boxed{\lambda_n^*=\lambda_d^*=
 \frac{h(m_n+m_d)}
 {m_d\sqrt{2m_nT}}\simeq8.6\,\mathrm{pm}}.
\]
\excheck 质心系总动量为0，
因此两粒子动量模相等，尽管质量不等。

\wavequantproblem{52}{238}
两相同非相对论粒子以相互垂直的方向运动，
在实验室系的德布罗意波长为
$\lambda_1,\lambda_2$。
求它们在质心系内各自的德布罗意波长。
\exsolution
质量均为 $m$，实验室动量
$\boldsymbol p_1\perp\boldsymbol p_2$，
$p_i=h/\lambda_i$。
质心速度
$\boldsymbol V=(\boldsymbol p_1+\boldsymbol p_2)/(2m)$。
第一粒子质心动量
$\boldsymbol p'_1=\boldsymbol p_1-m\boldsymbol V
=(\boldsymbol p_1-\boldsymbol p_2)/2$，
第二粒子为它的相反数，大小
$|\boldsymbol p'_1|=\tfrac12\sqrt{p_1^2+p_2^2}$。
于是
\[
 \boxed{\lambda^*=
 \frac{2}{\sqrt{1/\lambda_1^2+1/\lambda_2^2}}
 =\frac{2\lambda_1\lambda_2}
 {\sqrt{\lambda_1^2+\lambda_2^2}}}.
\]
\excheck 只有因为两实验室动量\emph{垂直}，
平方和才没有交叉项。

\wavequantproblem{53}{238}
室温下氢分子的\emph{最概然速率}对应
的德布罗意波长是多少？
\exsolution
麦克斯韦速度分布
$f(v)\propto v^2\exp[-mv^2/(2k_{\mathrm B}T)]$，
其速率峰值
$v_{\rm mp}=\sqrt{2k_{\mathrm B}T/m}$。
代入 $\lambda=h/(mv)$，
\[
 \boxed{\lambda(v_{\rm mp})=
 \frac{h}{\sqrt{2mk_{\mathrm B}T}}}.
\]
氢分子质量约 $m=2.016u$，
取 $T=300\,\mathrm K$，
\[
 \boxed{\lambda\simeq128\,\mathrm{pm}}.
\]
\excheck 本题问的是\emph{速率分布峰值对应的波长}，
与下一题“波长分布本身的峰值”不是同一概念。

\wavequantproblem{54}{238}
处于室温热平衡的氢分子，
若直接统计它们的德布罗意\emph{波长分布}，
求最概然波长。
\exsolution
设 $\lambda=h/(mv)$。
麦克斯韦速率分布
$f(v)\,dv\propto
v^2e^{-mv^2/(2k_{\mathrm B}T)}\,dv$。
进行概率变量替换：
\[
 \varphi(\lambda)\,d\lambda
 =f\!\left(\frac h{m\lambda}\right)
 \left|\frac{dv}{d\lambda}\right|d\lambda
 \propto\lambda^{-4}
 \exp\!\left[-\frac{h^2}
 {2mk_{\mathrm B}T\lambda^2}\right]d\lambda.
\]
极值条件 $d(\ln\varphi)/d\lambda=0$：
$-4/\lambda+h^2/(mk_{\mathrm B}T\lambda^3)=0$，
从而
\[
 \boxed{\lambda_{\rm mp}=
 \frac h{2\sqrt{mk_{\mathrm B}T}}
 \simeq0.089\,\mathrm{nm}\approx0.09\,\mathrm{nm}}.
\]
\excheck 概率密度在非线性变量变换下会改变形状；
不能把速率峰值代入 $\lambda=h/(mv)$ 当成波长峰值。

\wavequantproblem{55}{238}
推导动能为 $T$ 的相对论粒子的德布罗意波长；
若使用非相对论公式，要求相对误差不超过
$1\%$，分别求电子和质子的动能上限。
\exsolution
总能量 $E=mc^2+T$，
由相对论能量动量关系
$E^2=p^2c^2+m^2c^4$，得
\[
 pc=\sqrt{T(T+2mc^2)},\quad
 \boxed{\lambda_{\rm rel}=
 \frac{hc}{\sqrt{T(T+2mc^2)}}}.
\]
非相对论近似
$\lambda_{\rm nr}=h/\sqrt{2mT}$，
因此
\[
 \frac{\lambda_{\rm nr}}{\lambda_{\rm rel}}
 =\sqrt{1+\frac T{2mc^2}}.
\]
限定
$\lambda_{\rm nr}/\lambda_{\rm rel}-1\le0.01$，
有
\[
 \boxed{T\le2mc^2[(1.01)^2-1]
 =0.0402\,mc^2}.
\]
故
\[
 \boxed{T_{e,\max}\simeq20.5\,\mathrm{keV}},\qquad
 \boxed{T_{p,\max}\simeq37.7\,\mathrm{MeV}}.
\]
\excheck 1%是对\emph{波长}的相对误差；
非相对论波长比精确波长稍大。

\wavequantproblem{56}{238}
电子的德布罗意波长等于其普通
康普顿波长 $\lambda_C=h/(m_ec)$ 时，
其动能应为多少？
\exsolution
德布罗意式 $\lambda=h/p$，
令其等于 $h/(m_ec)$，得
$p=m_ec$。
由
$E=\sqrt{m_e^2c^4+p^2c^2}=\sqrt2m_ec^2$，
故动能
\[
 \boxed{T=(\sqrt2-1)m_ec^2
 \simeq0.212\,\mathrm{MeV}}.
\]
\excheck 所谓康普顿波长是 $h/(mc)$，
不是约化形式 $\hbar/(mc)$。

\wavequantproblem{57}{238}
X 射线管连续谱最短波长为
$\lambda_{\rm sh}=10.0\,\mathrm{pm}$。
求电子撞击阳极前的相对论德布罗意波长。
\exsolution
能量守恒给电子动能
$T=hc/\lambda_{\rm sh}$。
用相对论动量
$pc=\sqrt{T(T+2m_ec^2)}$，
于是
\[
 \boxed{\lambda_e=
 \frac{hc}{\sqrt{T(T+2m_ec^2)}}
 =\frac{\lambda_{\rm sh}}
 {\sqrt{1+\dfrac{2m_ec^2\lambda_{\rm sh}}{hc}}}}.
\]
数值 $T\simeq124\,\mathrm{keV}$，
\[
 \boxed{\lambda_e\simeq3.3\,\mathrm{pm}}.
\]
\excheck X 射线短波限是光子的波长，
不等于撞靶\emph{电子}的德布罗意波长。

\wavequantproblem{58}{238}
动能单一的电子束正入射宽
$b=1.0\,\mu\mathrm m$ 的窄方缝，
距缝 $\ell=50\,\mathrm{cm}$ 的屏上
单缝衍射中央主极大宽
$\Delta x=0.36\,\mathrm{mm}$。
求电子速度。
\exsolution
单缝第一暗纹满足
$b\sin\theta_1=\lambda_{\rm dB}$。
小角下 $\theta_1\simeq\lambda_{\rm dB}/b$，
中央主极大从左右第一暗纹之间的总宽度为
$\Delta x\simeq2\ell\lambda_{\rm dB}/b$，
即
\[
 \lambda_{\rm dB}=\frac{b\Delta x}{2\ell}.
\]
因 $\lambda_{\rm dB}=h/(m_ev)$，
\[
 \boxed{v=\frac{2h\ell}{m_eb\Delta x}
 \simeq2.0\times10^6\,\mathrm{m/s}}.
\]
\excheck 中央主极大宽度是从 $-\theta_1$ 到
$+\theta_1$ 的\emph{全宽}，不能少因子2。

\wavequantproblem{59}{238}
电子束经电压 $U=25\,\mathrm V$ 加速，
正入射两细缝，间距
$d=50\,\mu\mathrm m$。
屏距 $\ell=100\,\mathrm{cm}$，
求干涉条纹相邻亮纹距离。
\exsolution
加速电子获得动能
$T=eU$，非相对论动量
$p=\sqrt{2m_eeU}$，故波长
$\lambda=h/\sqrt{2m_eeU}$。
双缝相邻干涉亮纹角距满足
$\Delta\theta\simeq\lambda/d$，
屏上线距
\[
 \boxed{\Delta x=
 \frac{\ell h}{d\sqrt{2m_eeU}}
 \simeq4.9\,\mu\mathrm m}.
\]
\excheck 双缝条纹间距与双缝间距成反比，
且随加速电压按 $U^{-1/2}$ 缩小。

\wavequantproblem{60}{239}
动能单一的电子束以
$\theta=30^\circ$ 掠入射铝单晶天然晶面。
相邻晶面距 $d=0.20\,\mathrm{nm}$，
反射极大在某加速电压 $U_0$ 时出现。
若电压增为原来的
$\eta=2.25$ 倍，又出现相邻下一级反射极大，
求 $U_0$。
\exsolution
布拉格反射条件 $2d\sin\theta=m\lambda_m$。
非相对论电子
$\lambda\propto U^{-1/2}$，
因此电压变为 $\eta U_0$ 时
$\lambda_{\rm new}=\lambda_0/\sqrt\eta$。
对同一晶面和入射角，
$\lambda\propto1/m$，相邻两个整数级次满足
\[
 \frac{m+1}{m}=\sqrt\eta=1.5,
 \quad\Longrightarrow\quad m=2.
\]
所以
$2d\sin\theta=2\lambda_0$，
即 $\lambda_0=d\sin\theta$。
于是
\[
 \boxed{U_0=\frac{h^2}
 {2m_e e\,d^2\sin^2\theta}
 \simeq0.15\,\mathrm{kV}=150\,\mathrm V}.
\]
\excheck 电压提高使波长变短，
故对应的布拉格级次从2增到3。

\wavequantproblem{61}{239}
一束动能 $T=180\,\mathrm{eV}$ 的
单色电子正入射镍单晶，
在与晶体表面法线成
$\theta=55^\circ$ 的方向观察到\emph{第四级}
衍射反射极大。求晶面间距 $d$。
\exsolution
电子德布罗意波长
$\lambda=h/\sqrt{2m_eT}$。
入射方向沿向内的晶面法线、反射方向相对向外法线
成 $\theta$，两束传播方向夹角是 $\pi-\theta$；
布拉格掠射角为
$\alpha=(\pi-\theta)/2$，从而
$\sin\alpha=\cos(\theta/2)$。
第四级布拉格条件
$2d\sin\alpha=4\lambda$，得
\[
 \boxed{d=\frac{2\lambda}{\cos(\theta/2)}
 =\frac{2h}{\sqrt{2m_eT}\cos(\theta/2)}
 \simeq0.21\,\mathrm{nm}}.
\]
\excheck 若误将 $55^\circ$ 直接当作对晶面的
掠射角，会产生不同的结果；此题给的是相对于\emph{法线}的出射角。

\wavequantproblem{62}{239}
动能为 $T=10\,\mathrm{keV}$ 的电子束穿过
多晶铝箔，在距箔 $\ell=10.0\,\mathrm{cm}$
的屏上产生直径 $D=3.20\,\mathrm{cm}$ 的衍射环。
若该环由\emph{三级}布拉格反射形成，
求对应晶面间距。
\exsolution
设布拉格掠射角为 $\alpha$，
衍射圆锥半角为 $2\alpha$。
屏上环半径 $D/2$，所以
\[
 \tan(2\alpha)=\frac{D}{2\ell},
 \quad\alpha=\frac12\arctan\frac{D}{2\ell}
 \simeq4.55^\circ.
\]
电子波长在 $10\,\mathrm{keV}$ 下
$\lambda\simeq h/\sqrt{2m_eT}\simeq12.3\,\mathrm{pm}$
（精确相对论值约12.2pm）。
三级布拉格关系
$2d\sin\alpha=3\lambda$，
因此
\[
 \boxed{d=\frac{3\lambda}{2\sin\alpha}
 \simeq0.23\,\mathrm{nm}}.
\]
\excheck 圆环\emph{直径}先除2求半径；
衍射方向的偏转角为 $2\alpha$，不是 $\alpha$。

\wavequantproblem{63}{239}
电子被外部电压 $U$ 加速后入射金属，
金属内势（使电子动能增加）为
$U_i=15\,\mathrm V$。
求：(a) $U=150\,\mathrm V$ 时
电子波折射率 $n$；
(b) 电子波折射率偏离1不超过
$\eta=1.0\%$ 时 $U/U_i$ 应满足的范围。
\exsolution
金属外电子动能 $T_0=eU$，
金属内的动能是 $T_i=e(U+U_i)$。
动量随 $\sqrt T$ 增大，
故德布罗意波长随 $\sqrt T$ 减小，
定义折射率
\[
 \boxed{n=\frac{\lambda_{\rm out}}{\lambda_{\rm in}}
 =\sqrt{1+\frac{U_i}{U}}}.
\]
(a) 代数值
\[
 \boxed{n=\sqrt{1.10}\simeq1.049\approx1.05}.
\]
(b) 条件 $n-1\le\eta=0.010$ 给
$1+U_i/U\le(1+\eta)^2$，故
\[
 \boxed{\frac U{U_i}\ge
 \frac1{2\eta+\eta^2}\simeq49.75\approx50}.
\]
\excheck 内势为正增加电子动能，
所以 $n>1$；当 $U\to\infty$，$n\to1$。

\wavequantproblem{64}{239}
一维无限深方势阱宽度为 $\ell$。
若仅允许整数个德布罗意\emph{半波长}
恰好装入方阱，求粒子能级。
\exsolution
边界节点条件
$\ell=n\lambda/2$，
$n=1,2,\dots$。
因此
$\lambda=2\ell/n$，动量
$p=h/\lambda=nh/(2\ell)=n\pi\hbar/\ell$。
阱内势能为0，
\[
 \boxed{E_n=\frac{p^2}{2m}
 =\frac{n^2\pi^2\hbar^2}{2m\ell^2}},
 \quad n=1,2,\dots.
\]
\excheck $n=0$ 不对应可归一化的非零定态；
基态能量严格大于0。

\wavequantproblem{65}{239}
从德布罗意波的驻波观点解释玻尔角动量量子化：
证明氢原子电子只有沿周长能容纳
\emph{整数个完整德布罗意波长}的圆轨道才满足驻波条件。
\exsolution
电子沿半径 $r$ 的圆周运动，
动量为 $p=m_ev$，
德布罗意波长为
$\lambda=h/p$。
为使波函数沿圆轨道绕行一周后具有相同相位，
要求
\[
 2\pi r=n\lambda,\qquad n=1,2,\dots.
\]
代入 $\lambda=h/p$，
\[
 2\pi rp=nh,\quad
 \boxed{L=rp=n\hbar}.
\]
这恰为玻尔的圆轨道角动量量子化条件。
\excheck 非整数圈数会导致同一点相位不连续，
无法形成稳定单值驻波；
此结论是旧量子理论的波动解释，非完整现代氢原子轨道理论。
% === END 35-wave-properties-049-065.tex ===

% === BEGIN 35-wave-properties-066-080.tex blob 8902ae626aa2ea6114a37c6cf7dd4708e06a0d2f ===
% Irodov Part VI §6.2, original 66--80. Printed pages 239--240.
% User-provided Arihant PDF pages 248--249; reference answers PDF p.350.
% Uncertainty estimates explicitly distinguish hbar and Delta x conventions.
\wavequantproblem{66}{239}
已知电子、质子以及质量为 $m_b=1\,\mathrm{mg}$ 的小球，其位置（小球为质心位置）的不确定度均为 $\Delta x=1\,\mu\mathrm m$。估算测量其速度时的最小不确定度。
\exsolution
海森堡不确定关系
$\Delta x\,\Delta p_x\geq\hbar/2$，
对低速物体 $\Delta p_x=m\Delta v_x$，故
\[
 \Delta v_{x,\min}=\frac{\hbar}{2m\Delta x}.
\]
取 $m_e=9.109\times10^{-31}\,\mathrm{kg}$、
$m_p=1.673\times10^{-27}\,\mathrm{kg}$、
$m_b=10^{-6}\,\mathrm{kg}$，得
\[
 \boxed{\Delta v_{e,\min}=5.79\times10^1\,\mathrm{m/s}},
 \quad\boxed{\Delta v_{p,\min}=3.15\times10^{-2}\,\mathrm{m/s}},
 \quad\boxed{\Delta v_{b,\min}=5.27\times10^{-23}\,\mathrm{m/s}}.
\]
\excheck 若原书按数量级用 $\Delta x\Delta p\sim\hbar$，
数值可差一个2的系数；三者的速度精度差异随质量倒数变化，
这不影响微观与宏观比较结论。

\wavequantproblem{67}{239}
利用不确定原理估算电子在氢原子中的速度不确定度，取原子线尺度 $l=0.10\,\mathrm{nm}$，并与第一玻尔轨道的电子速度比较。
\exsolution
当电子的位置被限制在 $l$ 的量级，
动量不确定度在数量级上为
$\Delta p\sim\hbar/l$，所以
\[
 \boxed{\Delta v\sim\frac{\hbar}{m_e l}
 \simeq1.16\times10^6\,\mathrm{m/s}}.
\]
第一玻尔轨道电子速度
$v_1=\alpha_{\rm fs}c\simeq2.19\times10^6\,\mathrm{m/s}$，
从而
\[
 \boxed{\frac{\Delta v}{v_1}\sim0.53}.
\]
\excheck 本题为数量级估算，而非通过精确几率分布求标准差；
电子速度的不确定量确实与其典型轨道速度同阶。

\wavequantproblem{68}{239}
证明当粒子的坐标不确定度为 $\Delta x=\lambda/(2\pi)$，
其中 $\lambda$ 为其德布罗意波长时，
它的速度不确定度与粒子速度本身同数量级。
\exsolution
由 $\lambda=h/p=2\pi\hbar/p$ 得
$\Delta x=\hbar/p$。
不确定性关系给出
\[
 \Delta p\geq\frac{\hbar}{2\Delta x}=\frac p2.
\]
若粒子处于非相对论区，
$p=mv$、$\Delta p=m\Delta v$，
于是
\[
 \boxed{\Delta v\geq\frac v2}.
\]
\excheck 系数 $1/2$ 依具体标准差定义；
结论是$\Delta v$不能比$v$小很多，即为同阶。

\wavequantproblem{69}{240}
一个自由电子最初局限在长度量级
$l=0.10\,\mathrm{nm}$ 的空间区域，
用不确定原理估计其波包宽度增至原来
$\eta=10$ 倍所需时间。
\exsolution
初始局域化使动量展宽在数量级上
$\Delta p\sim\hbar/l$，
速度展宽为 $\Delta v\sim\hbar/(m_e l)$。
自由传播中不同速度的波包分量逐渐分离，
波包额外扩展距离约 $\Delta v\,t$。
当展宽达到 $\eta l$ 的量级时
\[
 \eta l\sim\frac{\hbar t}{m_e l},
 \qquad
 \boxed{t\sim\frac{\eta m_e l^2}{\hbar}
 \simeq8.6\times10^{-16}\,\mathrm s}.
\]
\excheck 这是数量级估算，精确高斯波包的展开时间还取决于
$l$ 是振幅宽度、概率标准差还是其他宽度约定；
不宜将题中的“估计”误写成唯一精确时间。

\wavequantproblem{70}{240}
利用不确定原理估算局域于
$l=0.20\,\mathrm{nm}$ 的电子的最小动能。
\exsolution
电子被束缚在长度$l$量级时，
$\Delta p\sim\hbar/l$，
取典型动量$p\sim\Delta p$，
\[
 \boxed{T_{\min}\sim\frac{\hbar^2}{2m_e l^2}
 \simeq0.95\,\mathrm{eV}\approx1\,\mathrm{eV}}.
\]
\excheck 上式是原书要求的粗估下限量级；
若指定精确无限深一维方阱边界，
其基态能级则为$\pi^2\hbar^2/(2m_el^2)$，
有额外的$\pi^2$系数。

\wavequantproblem{71}{240}
电子动能约 $T=4\,\mathrm{eV}$，
被局限在尺寸$l=1\,\mu\mathrm m$的区域。
估算相对速度不确定度。
\exsolution
由不确定关系按量级
$\Delta p\sim\hbar/l$；
非相对论动量$p=\sqrt{2m_eT}$，
所以
\[
 \boxed{\frac{\Delta v}{v}\sim\frac{\Delta p}{p}
 \sim\frac{\hbar}{l\sqrt{2m_eT}}
 \simeq9.8\times10^{-5}\approx10^{-4}}.
\]
\excheck 由于$l$远大于电子的德布罗意波长，
相对速度展宽远小于1，量级合理。

\wavequantproblem{72}{240}
电子位于宽度为$l$的无限深一维方阱，
用不确定原理估算占据最低可能能量时
电子施加在阱壁上的力。
\exsolution
最小动量量级为
$p\sim\hbar/l$，
相应能量为
$E_{\min}\sim\hbar^2/(2m_el^2)$。
缓慢改变壁间距所需的功
$dW=F\,dl$，
故电子对可移动阱壁的压力效应
\[
 \boxed{F\sim-\frac{dE_{\min}}{dl}
 \sim\frac{\hbar^2}{m_el^3}}.
\]
\excheck 精确无限阱基态是
$E_1=\pi^2\hbar^2/(2m_el^2)$，
对应的精确广义力为
$F_1=\pi^2\hbar^2/(m_el^3)$；
原书要求的是不确定原理估算。

\wavequantproblem{73}{240}
质量为$m$的粒子在一维简谐势
$U(x)=kx^2/2$中运动。
用不确定原理估计最低允许能量，
并与精确解比较。
\exsolution
若位置展宽为$x$的量级，
则$p\sim\hbar/x$，因而估计总能量
\[
 E(x)\sim\frac{\hbar^2}{2mx^2}+\frac{kx^2}{2}.
\]
极小值条件
$dE/dx=-\hbar^2/(mx^3)+kx=0$，
给出
$x^4=\hbar^2/(mk)$，
于是
\[
 \boxed{E_{\min}\sim\hbar\sqrt{k/m}
 =\hbar\omega},\qquad \omega=\sqrt{k/m}.
\]
\excheck 采用真正的标准差不确定关系与高斯态变分
可以得到精确基态能
$E_0=\hbar\omega/2$；旧式量级估算只保证尺度。

\wavequantproblem{74}{240}
用不确定原理估计氢原子电子所允许的
最低能量及对应的最概然量级核间距离。
\exsolution
电子局域到尺度$r$，按量级
$p\sim\hbar/r$；总能量近似
\[
 E(r)\sim\frac{\hbar^2}{2m_er^2}
 -\frac{k_ee^2}{r},
 \qquad k_e=\frac1{4\pi\varepsilon_0}.
\]
令$dE/dr=0$，
$-\hbar^2/(m_er^3)+k_ee^2/r^2=0$，得
\[
 \boxed{r_*\sim\frac{\hbar^2}{m_ek_ee^2}
 =a_0\simeq52.9\,\mathrm{pm}}.
\]
代回能量
\[
 \boxed{E_{\min}\sim-\frac{m_e(k_ee^2)^2}{2\hbar^2}
 \simeq-13.6\,\mathrm{eV}}.
\]
\excheck 用等式求极小值只是在估算有效尺度；
真正三维氢基态径向概率最大点恰为$a_0$，
但这需用量子波函数另行证明。

\wavequantproblem{75}{240}
速度$v=600\,\mathrm{m/s}$的平行氢原子束
垂直照射狭缝，缝后距离$\ell=1.0\,\mathrm m$
处放置屏。用不确定原理求成像总宽度最小时
狭缝宽度$\delta$。
\exsolution
狭缝横向局域化给
$\Delta p_y\sim\hbar/\delta$。
原子纵向动量$p_x\simeq m_Hv$，
角展宽$\Delta\theta\sim\Delta p_y/p_x$，
传播距离$\ell$引起的附加宽度为
$\Delta_{\rm dif}\sim\ell\hbar/(m_Hv\delta)$。
总宽度估算
\[
 W(\delta)\sim\delta+\frac{\hbar\ell}{m_Hv\,\delta}.
\]
对$\delta$求极值，得
\[
 \boxed{\delta_{\min}\sim\sqrt{\frac{\hbar\ell}{m_Hv}}
 \simeq1.03\times10^{-5}\,\mathrm m
 \approx0.01\,\mathrm{mm}}.
\]
\excheck 增大缝宽会直接增加几何模糊，
减小缝宽则会增加衍射模糊，因此存在非零最优值。

\wavequantproblem{76}{240}
求一质量为$m$的自由粒子的
含时薛定谔方程的一个特解。
\exsolution
自由粒子$U=0$，满足
\[
 i\hbar\frac{\partial\Psi}{\partial t}
 =-\frac{\hbar^2}{2m}\nabla^2\Psi.
\]
取平面波试探解
$\Psi=A\exp[i(\boldsymbol k\cdot\boldsymbol r-\omega t)]$。
代入得
$\hbar\omega=\hbar^2k^2/(2m)$，
所以一个特解是
\[
 \boxed{\Psi(\boldsymbol r,t)=
 A\exp\!\left[\frac{i}{\hbar}
 \left(\boldsymbol p\cdot\boldsymbol r-\frac{p^2}{2m}t\right)\right]},
 \quad\boldsymbol p=\hbar\boldsymbol k.
\]
\excheck 无限空间中的单个平面波不是平方可积函数，
描述可归一化波包时须作动量态的叠加。

\wavequantproblem{77}{240}
宽度为$l$的无限深一维势阱$0<x<l$
内粒子处于基态。求它位于
$l/3\le x\le2l/3$的概率。
\exsolution
归一化基态波函数
$\psi_1(x)=\sqrt{2/l}\sin(\pi x/l)$。
因此
\[
 P=\frac2l\int_{l/3}^{2l/3}\sin^2\frac{\pi x}{l}\,dx.
\]
利用$\sin^2u=(1-\cos2u)/2$，
\[
 \boxed{P=\frac13+\frac{\sqrt3}{2\pi}
 \simeq0.609}.
\]
\excheck 概率超过几何区间比例$1/3$，
因为基态概率密度在阱中部最大。

\wavequantproblem{78}{240}
无限深一维方势阱以阱中心为坐标原点，
宽度为$l$，即$|x|<l/2$。
求各定态的归一化波函数。
\exsolution
阱内定态方程
$\psi''+k^2\psi=0$；
边界$\psi(\pm l/2)=0$，
故$k=n\pi/l$、$n=1,2,\dots$。
按奇偶性写成
\[
 \boxed{\psi_n(x)=
 \begin{cases}
 \sqrt{2/l}\cos(n\pi x/l),&n=1,3,5,\ldots,\\
 \sqrt{2/l}\sin(n\pi x/l),&n=2,4,6,\ldots,
 \end{cases}}
\quad |x|<l/2,
\]
阱外$\psi_n=0$。
对应能量
$E_n=n^2\pi^2\hbar^2/(2ml^2)$。
\excheck 余弦/正弦只差可忽略的整体符号；
归一化需满足$\int_{-l/2}^{l/2}|\psi_n|^2dx=1$。

\wavequantproblem{79}{240}
证明无限深一维方阱$0<x<l$
内不同能级的定态波函数彼此正交。
\exsolution
取标准归一化本征函数
$\psi_n(x)=\sqrt{2/l}\sin(n\pi x/l)$，
$n=1,2,\dots$。
对$n\ne n'$，用积化和差：
\[
 \int_0^l\psi_n\psi_{n'}dx
 =\frac1l\int_0^l
 \left[\cos\frac{(n-n')\pi x}{l}
 -\cos\frac{(n+n')\pi x}{l}\right]dx=0.
\]
因此
\[
 \boxed{\langle n|n'\rangle=\delta_{nn'}}.
\]
\excheck $n=n'$时积分为1；
正交性可独立由厄米哈密顿算符不同本征值证明。

\wavequantproblem{80}{240}
电子位于长$l$的无限深一维势阱，
能级极为密集。推导单位能量间隔内能级数
$dN/dE$，并求$l=1.0\,\mathrm{cm}$、
$E=1.0\,\mathrm{eV}$时的值。
\exsolution
精确能级
$E_n=n^2\pi^2\hbar^2/(2m_el^2)$。
高$n$连续近似下
$n(E)=l\sqrt{2m_eE}/(\pi\hbar)$。
对$E$求导，得
\[
 \boxed{\frac{dN}{dE}=\frac l{\pi\hbar}
 \sqrt{\frac{m_e}{2E}}}.
\]
此式以J为能量单位。若改为每eV的态数，
须乘$1\,\mathrm{eV}=1.602176634\times10^{-19}\,\mathrm J$，
故
\[
 \boxed{\left.\frac{dN}{dE}\right|_{1\,\mathrm{eV}}
 \simeq8.1\times10^6\ \mathrm{eV^{-1}}}.
\]
\excheck 连续密度只适用于$E_n$相邻间隔远小于所考察能量范围，
且这里未计电子自旋二重简并。
% === END 35-wave-properties-066-080.tex ===

% === BEGIN 35-wave-properties-081-096.tex blob d0be10d8e7fadc5e264ee5d78dd631234c85f7be ===
% Irodov Part VI §6.2 original problems 81--96, closing §6.2.
% Original printed pp.241--243, user PDF pages 250--252;
% answer pp.342--344, user PDF pages 351--353.
% Scattering and barrier diagram conventions read from original figures 6.2--6.6.
\wavequantproblem{81}{241}
质量为$m$的粒子被局限在边长为$l_1,l_2$的二维
无限深矩形势阱中。(a) 求全部允许能量；
(b) 若势阱为边长$l$的正方形，求前四个不同能级。
\exsolution
势阱内的定态薛定谔方程可分离：
$\psi(x,y)=X(x)Y(y)$。
边界上波函数为零，故
$X\propto\sin(n_1\pi x/l_1)$、
$Y\propto\sin(n_2\pi y/l_2)$，
$n_1,n_2=1,2,\ldots$。得到
\[
 \boxed{E_{n_1n_2}=\frac{\pi^2\hbar^2}{2m}
 \left(\frac{n_1^2}{l_1^2}+\frac{n_2^2}{l_2^2}\right)}.
\]
正方形时令
$\varepsilon=\pi^2\hbar^2/(2ml^2)$；
最低四个\emph{不同}的$n_1^2+n_2^2$为
$2,5,8,10$，因此
\[
 \boxed{E_1=2\varepsilon,\quad E_2=5\varepsilon,\quad
 E_3=8\varepsilon,\quad E_4=10\varepsilon}.
\]
\excheck 二维正方阱中$(n_1,n_2)$与$(n_2,n_1)$
常产生简并；“第四能级”应按不同能量而非重复状态排序。

\wavequantproblem{82}{241}
粒子在二维无限深矩形阱
$0<x<a,\ 0<y<b$中处于基态，
求其位于$0<x<a/3$区域的概率。
\exsolution
归一化基态波函数
$\psi_{11}=\sqrt{4/(ab)}
\sin(\pi x/a)\sin(\pi y/b)$。
对全部$y$积分为1，
故
\[
 P=\frac2a\int_0^{a/3}\sin^2\frac{\pi x}{a}\,dx
 =\boxed{\frac13-\frac{\sqrt3}{4\pi}
 \simeq0.1955}.
\]
\excheck 由于基态在$x=0$的概率密度为0，
左侧三分之一的概率比几何面积份额$1/3$小。

\wavequantproblem{83}{241}
质量$m$的粒子处于边长$a$的三维无限深立方势阱。
求：(a)全部本征能量；
(b)第三、第四不同能级的间隔；
(c)第六个不同能级的能量及简并度。
\exsolution
沿三坐标可分离并施加端点节点条件，
\[
 \boxed{E_{n_1n_2n_3}=
 \frac{\pi^2\hbar^2}{2ma^2}
 (n_1^2+n_2^2+n_3^2)},\quad n_j\ge1.
\]
最低六个\emph{不同}的平方和依次为
$3,6,9,11,12,14$，例如第三能级属于$(1,2,2)$，
第四能级属于$(1,1,3)$。
记$\varepsilon=\pi^2\hbar^2/(2ma^2)$，
所以
\[
 \boxed{E_4-E_3=(11-9)\varepsilon
 =\frac{\pi^2\hbar^2}{ma^2}}.
\]
第六能级对应$(1,2,3)$的六种排列，
\[
 \boxed{E_6=14\varepsilon=\frac{7\pi^2\hbar^2}{ma^2}},
 \qquad \boxed{g_6=3!=6}.
\]
\excheck 简并度不含自旋；
若额外计自旋自由度，需要另外乘相应因子。

\wavequantproblem{84}{241}
从定态薛定谔方程证明：
粒子势能$U(x)$在某点发生\emph{有限跃变}时，
波函数一阶导数仍连续。
\exsolution
定态方程
\[
 -\frac{\hbar^2}{2m}\psi''(x)+U(x)\psi(x)=E\psi(x).
\]
将方程在跃变点$x_0$两侧
$[x_0-\delta,x_0+\delta]$上积分，
\[
 \psi'(x_0+\delta)-\psi'(x_0-\delta)
 =\frac{2m}{\hbar^2}
 \int_{x_0-\delta}^{x_0+\delta}(U-E)\psi\,dx.
\]
若$U$有界且$\psi$有限，
右边随$\delta\to0$趋于零，
因此
\[
 \boxed{\psi'(x_0^+)=\psi'(x_0^-)}.
\]
\excheck 有限跃变也要求$\psi$连续；
$\delta$型奇异势不属于本题，
其$\psi'$可出现有限跳跃。

\wavequantproblem{85}{241}
粒子位于图6.2所示的单侧有限深方势阱：
$x\le0$处为无限高势壁，
$0<x<l$有$U=0$，
$x\ge l$有$U=U_0>0$。
(a) 推导$E<U_0$的束缚能级方程并表明能谱离散；
(b) 求存在第$n$个束缚态所需$l^2U_0$的阈值。
\exsolution
令
$k=\sqrt{2mE}/\hbar$、
$\kappa=\sqrt{2m(U_0-E)}/\hbar$。
归一化束缚解在阱内、阱外分别是
$\psi_<=A\sin(kx)$、
$\psi_>=B e^{-\kappa(x-l)}$。
在$x=l$令波函数及一阶导数连续，
\[
 A\sin(kl)=B,\quad
 Ak\cos(kl)=-B\kappa,
\]
于是
\[
 \boxed{k\cot(kl)=-\kappa
 \quad\Longleftrightarrow\quad \tan(kl)=-\frac k\kappa}.
\]
又$k^2+\kappa^2=2mU_0/\hbar^2$，
\[
 \boxed{\sin(kl)=\pm\frac{\hbar k}{\sqrt{2mU_0}}},
\]
同时必须选取$\tan(kl)<0$的根。
离散实根给出离散本征能级。
在第$n$个态刚好出现时$E\uparrow U_0$、
$\kappa\downarrow0$，故
$kl=(n-\tfrac12)\pi$，从而
\[
 \boxed{(l^2U_0)_{n,\mathrm{th}}=
 \frac{(2n-1)^2\pi^2\hbar^2}{8m}},
 \quad n=1,2,\dots.
\]
\excheck 阈值处$E=U_0$的态已不严格平方可积，
因此真正束缚态须满足严格大于阈值。
高斯/文本抽取常漏去$(2n-1)^2$的平方。

\wavequantproblem{86}{241}
沿用前题单侧有限深阱。
当$E=U_0/2$且
$l^2U_0=(3\pi/4)^2\hbar^2/m$时，
求粒子出现在$x>l$区域的概率。
\exsolution
此时
$k=\kappa=\sqrt{mU_0}/\hbar$，
$kl=3\pi/4$，
满足$\tan(kl)=-1$。
取
\[
 \psi(x)=
 \begin{cases}
 A\sin(kx),&0<x<l,\\
 A\sin(kl)e^{-k(x-l)},&x>l.
 \end{cases}
\]
阱外积分
$I_{\rm out}=|A|^2\sin^2(kl)/(2k)=|A|^2/(4k)$；
阱内积分
\[
 I_{\rm in}=|A|^2
 \left[\frac l2-\frac{\sin(2kl)}{4k}\right]
 =\frac{|A|^2}{4k}\left(1+\frac{3\pi}{2}\right).
\]
所以
\[
 \boxed{P_{\rm out}=
 \frac{I_{\rm out}}{I_{\rm out}+I_{\rm in}}
 =\frac2{4+3\pi}\simeq0.149}.
\]
\excheck 类禁区概率不为零是波函数指数尾部的结果，
不是经典粒子穿透并持续向右运动的定态概率流。

\wavequantproblem{87}{242}
质量$m$的粒子在半径$r_0$的三维球形
无限深势阱中运动：
$U=0$（$r<r_0$）、
$U=\infty$（$r\ge r_0$）。
若波函数仅与径向$r$有关，求允许能量。
\exsolution
径向定态方程（角动量$l_{\rm orb}=0$）为
\[
 \psi''+\frac2r\psi'+k^2\psi=0,\quad
 k^2=\frac{2mE}{\hbar^2}.
\]
令$\chi(r)=r\psi(r)$，则
$\chi''+k^2\chi=0$。
原点波函数有限要求$\chi(0)=0$，
球壁要求$\chi(r_0)=0$，
因此
$\chi\propto\sin(kr)$且$kr_0=n\pi$。
所以
\[
 \boxed{E_n=\frac{n^2\pi^2\hbar^2}{2mr_0^2}},
 \quad n=1,2,\dots.
\]
\excheck 这是只考虑球对称$s$波的能谱，
并非所有角动量分波的完整三维球阱谱。

\wavequantproblem{88}{242}
在前题球形无限深势阱中：
(a) 求只依赖$r$的归一化波函数；
(b) 求基态最概然径向$r_{\rm mp}$以及
$P(r<r_{\rm mp})$。
\exsolution
设$\psi_n=A_n\sin(n\pi r/r_0)/r$。
归一化积分
\[
 1=4\pi|A_n|^2\int_0^{r_0}
 \sin^2\frac{n\pi r}{r_0}\,dr
 =2\pi r_0|A_n|^2,
\]
故
\[
 \boxed{\psi_n(r)=\frac1{\sqrt{2\pi r_0}}
 \frac{\sin(n\pi r/r_0)}r},\quad r<r_0.
\]
基态的径向概率密度
$dP/dr=4\pi r^2|\psi_1|^2
=(2/r_0)\sin^2(\pi r/r_0)$，
在$r=r_0/2$达最大。因此
\[
 \boxed{r_{\rm mp}=\frac{r_0}2},\qquad
 \boxed{P(r<r_{\rm mp})=\frac2{r_0}
 \int_0^{r_0/2}\sin^2\frac{\pi r}{r_0}dr=\frac12}.
\]
\excheck 最概然\emph{径向}应最大化$4\pi r^2|\psi|^2$，
不是单独最大化$|\psi(r)|^2$。

\wavequantproblem{89}{242}
粒子在球形有限深势阱中：
$U(r)=0$（$r<r_0$）、
$U(r)=U_0$（$r\ge r_0$）。
(a) 对球对称定态导出$E<U_0$束缚态能级条件；
(b) 求出现第一个束缚态所需$r_0^2U_0$阈值。
\exsolution
仍令$\chi=r\psi$，
则径向方程化为一维形式。
阱内
$\chi_<=A\sin(kr)$，
阱外取衰减解
$\chi_>=B e^{-\kappa(r-r_0)}$，
其中$k=\sqrt{2mE}/\hbar$、
$\kappa=\sqrt{2m(U_0-E)}/\hbar$。
在$r=r_0$匹配$\chi,\chi'$，
\[
 \boxed{k\cot(kr_0)=-\kappa},
 \quad
 \boxed{\sin(kr_0)=\pm\frac{\hbar k}{\sqrt{2mU_0}}}.
\]
第一次出现束缚态的阈值在
$\kappa\to0$、$kr_0\to\pi/2$，
\[
 \boxed{(r_0^2U_0)_{\mathrm{th}}=
 \frac{\pi^2\hbar^2}{8m}}.
\]
\excheck 对三维\emph{球形}有限势阱必须超过一定深度
才支持$s$波束缚态；
这一点不同于对称的一维有限势阱。

\wavequantproblem{90}{242}
质量为$m$的粒子处于一维简谐势
$U(x)=kx^2/2$。
已知其基态波函数形式为
$\psi(x)=A e^{-\alpha x^2}$（$\alpha>0$）。
从薛定谔方程求$\alpha$与基态能量。
\exsolution
计算导数
$\psi''=(4\alpha^2x^2-2\alpha)\psi$。
代入
\[
 -\frac{\hbar^2}{2m}\psi''+\frac{kx^2}{2}\psi=E\psi
\]
得
\[
 \left[\frac{\hbar^2\alpha}{m}
 +\left(\frac k2-\frac{2\hbar^2\alpha^2}{m}\right)x^2\right]\psi
 =E\psi.
\]
因等式对一切$x$成立，$x^2$系数为0，
$\alpha=\sqrt{mk}/(2\hbar)$，
而常数项即能量。故
\[
 \boxed{\alpha=\frac{m\omega}{2\hbar}},\quad
 \boxed{E_0=\frac{\hbar\omega}{2}},
 \qquad \omega=\sqrt{k/m}.
\]
\excheck 这是严格量子基态，比粗略不确定度估算题73
所得$\sim\hbar\omega$多出明确的$1/2$系数。

\wavequantproblem{91}{242}
氢原子电子某一球对称定态的波函数为
$\psi(r)=A(1+ar)e^{-\alpha r}$，其中
$A,a,\alpha$均为常数。求其能量及主量子数。
\exsolution
氢原子库仑势
$U=-k_ee^2/r$。
已知球对称态的径向薛定谔方程
\[
 -\frac{\hbar^2}{2m_e}
 \left(\psi''+\frac2r\psi'\right)
 -\frac{k_ee^2}{r}\psi=E\psi.
\]
令$f=(1+ar)e^{-\alpha r}$，
则
\[
 f'=e^{-\alpha r}(a-\alpha-a\alpha r),\quad
 f''=e^{-\alpha r}(\alpha^2-2a\alpha+a\alpha^2r).
\]
代入并分别比较常数、$r^{-1}$及$r$项，
求得
$\alpha=1/(2a_0)$、$a=-1/(2a_0)$，
其中$a_0=\hbar^2/(m_ek_ee^2)$。
于是该态正是氢原子的$2s$径向态，
\[
 \boxed{E=-\frac{m_e(k_ee^2)^2}{8\hbar^2}
 =-\frac{13.6}{4}\,\mathrm{eV}=-3.40\,\mathrm{eV}},
 \qquad\boxed{n=2}.
\]
\excheck 虽然波函数只有$r$，仍可有$n=2$的径向节点，
不能由角动量为0就断定$n=1$。

\wavequantproblem{92}{242}
基态氢原子电子波函数
$\psi=Ae^{-r/a_0}$，
$a_0$为第一玻尔半径。
求：(a)最概然核间距离；
(b)电子所受库仑力模的平均值；
(c)平均势能。
\exsolution
归一化给
$A=1/\sqrt{\pi a_0^3}$。
径向概率密度
$f(r)=4r^2e^{-2r/a_0}/a_0^3$，
令$f'(r)=0$得
\[
 \boxed{r_{\rm mp}=a_0}.
\]
库仑力模为$F(r)=k_ee^2/r^2$，
所以
\[
 \langle F\rangle=
 \int_0^\infty f(r)\frac{k_ee^2}{r^2}\,dr
 =\boxed{\frac{2k_ee^2}{a_0^2}}.
\]
势能$U(r)=-k_ee^2/r$，则
\[
 \boxed{\langle U\rangle=
 -\frac{k_ee^2}{a_0}}.
\]
\excheck 力模的\emph{平均值}不是
$k_ee^2/\langle r\rangle^2$；
必须先按波函数求概率平均。

\wavequantproblem{93}{242}
基态氢原子电子波函数
$\psi=Ae^{-r/a_0}$。
求在氢原子\emph{原子核中心}，
该电子电荷云产生的平均静电势。
\exsolution
电荷密度
$\rho(r)=-e|\psi(r)|^2$，
其中$|A|^2=1/(\pi a_0^3)$。
球对称分布在原点的势为
\[
 \varphi(0)=k_e\int_{\mathbb R^3}
 \frac{\rho(\boldsymbol r)}r\,d^3r
 =4\pi k_e\int_0^\infty \rho(r)\,r\,dr.
\]
代入积分
$\int_0^\infty r e^{-2r/a_0}\,dr=a_0^2/4$，
得
\[
 \boxed{\varphi(0)=-\frac{k_ee}{a_0}
 \simeq-27.2\,\mathrm V}.
\]
\excheck 此处仅计算\emph{电子云}在核位置的势，
不包括点核自身的发散自势。

\wavequantproblem{94}{243}
图6.3中势阶跃
$U(x)=0$（$x<0$），
$U(x)=U_0$（$x>0$）。
求：(a)当$E>U_0$时从左方入射粒子的反射系数；
(b)当$E<U_0$时，右侧禁区中概率密度降到界面处
$1/e$所需渗入深度。
\exsolution
(a) 两区域波数
$k_1=\sqrt{2mE}/\hbar$、
$k_2=\sqrt{2m(E-U_0)}/\hbar$。
写
$\psi_1=Ae^{ik_1x}+Be^{-ik_1x}$、
$\psi_2=Ce^{ik_2x}$。
在$x=0$匹配$\psi$与$\psi'$，
由$A+B=C$、
$k_1(A-B)=k_2C$，
\[
 \boxed{R=\frac{|B|^2}{|A|^2}
 =\left(\frac{k_1-k_2}{k_1+k_2}\right)^2}.
\]
(b) 若$E<U_0$，
$\psi_2=C e^{-\kappa x}$，
$\kappa=\sqrt{2m(U_0-E)}/\hbar$。
概率密度衰减为$e^{-2\kappa x}$，
满足$e^{-2\kappa x_{\rm eff}}=e^{-1}$，
故
\[
 \boxed{x_{\rm eff}=\frac1{2\kappa}
 =\frac{\hbar}{2\sqrt{2m(U_0-E)}}}.
\]
\excheck 本题“下降$e$倍”的对象是\emph{概率密度}，
不是波函数振幅；两者的衰减长度相差2倍。

\wavequantproblem{95}{243}
利用半经典WKB公式，求能量为$E<U_0$的电子穿过
图6.4中的矩形势垒和图6.5中的线性上升三角势垒的
透射概率指数近似。矩形宽$l$高$U_0$；
三角势垒取$U(x)=U_0(x+l)/l$于$-l<x<0$，
在$x=0$降至零。
\exsolution
WKB禁止区透射率的主指数是
\[
 D\simeq\exp\!\left[-\frac2\hbar
 \int_{x_-}^{x_+}\sqrt{2m(U(x)-E)}\,dx\right].
\]
(a) 矩形势垒禁止区贯穿整个宽度，
\[
 \boxed{D_{\rm rect}\simeq
 \exp\!\left[-\frac{2l}{\hbar}\sqrt{2m(U_0-E)}\right]}.
\]
(b) 三角势垒的左转折点
$x_t=-l+lE/U_0$。
由$dU=(U_0/l)\,dx$，
\[
 \int_{x_t}^{0}\sqrt{2m(U-E)}\,dx
 =\frac{l\sqrt{2m}}{U_0}
 \int_E^{U_0}\sqrt{U-E}\,dU
 =\frac{2l\sqrt{2m}}{3U_0}(U_0-E)^{3/2}.
\]
故按图6.5所示\emph{单侧}斜坡，
\[
 \boxed{D_{\rm tri}\simeq
 \exp\!\left[-\frac{4l\sqrt{2m}}{3\hbar U_0}
 (U_0-E)^{3/2}\right]}.
\]
\excheck 用户所附版本书末第95(b)题印有
$8l/(3\hbar U_0)$而非本式的$4l/(3\hbar U_0)$；
依原图仅一个宽度$l$的线性斜坡时，严格的WKB
积分是系数4；若将禁区积分路径重复两次才出现8。
该冲突应保留供二校，不能悄悄修改图或积分区间。

\wavequantproblem{96}{243}
用WKB公式求质量$m$、能量$E<U_0$的粒子
穿过图6.6所示抛物势垒
$U(x)=U_0(1-x^2/l^2)$的透射概率指数近似。
\exsolution
经典转折点由$U(x_\pm)=E$求出
$x_\pm=\pm a$，
$a=l\sqrt{1-E/U_0}$。
在禁止区
$U(x)-E=(U_0-E)(1-x^2/a^2)$，
因此
\[
 \int_{-a}^{a}\sqrt{2m(U-E)}\,dx
 =a\sqrt{2m(U_0-E)}
 \int_{-1}^{1}\sqrt{1-u^2}\,du.
\]
利用半圆面积积分
$\int_{-1}^{1}\sqrt{1-u^2}\,du=\pi/2$，
WKB指数为
\[
 \boxed{D\simeq
 \exp\!\left[-\frac{\pi l}{\hbar}
 \sqrt{\frac{2m}{U_0}}\,(U_0-E)\right]}.
\]
\excheck 势垒宽度不应一律取$2l$，
应当按能量$E$确定转折点$\pm a$；
所得指数与原书第96题答案相符。
% === END 35-wave-properties-081-096.tex ===
